% Sorgente LaTeX canonico, importato e revisionato dal precedente sorgente Markdown.
\chapter{Integrali impropri}
\label{chap:16}

\begin{obiettivocapitolo}{decidere convergenza e calcolare integrali con estremi infiniti o singolarità}
\prioritaalta. Ogni punto improprio va trattato con un limite separato.
\end{obiettivocapitolo}
\begin{proceduraesame}{integrale improprio}
\begin{enumerate}
\item Identifica tutti gli estremi infiniti e le singolarità interne o agli estremi.
\item Spezza l'integrale in modo che ciascun pezzo contenga una sola causa di improprietà.
\item Sostituisci l'estremo problematico con un parametro e scrivi il limite corretto.
\item Prima di calcolare, confronta localmente con integrali campione $x^{-p}$ o con equivalenti.
\item L'integrale converge soltanto se ogni limite è finito; altrimenti diverge.
\end{enumerate}
\end{proceduraesame}
\begin{sceltametodo}{criterio rapido}
All'infinito, $\int_1^\infty x^{-p}dx$ converge per $p>1$. Vicino a zero, $\int_0^1 x^{-p}dx$ converge per $p<1$. Per funzioni non negative usa confronto o confronto asintotico; per segno variabile verifica prima la convergenza assoluta.
\end{sceltametodo}
\begin{errorefrequente}{valore principale}
Il valore principale di Cauchy può esistere per compensazione simmetrica anche quando l'integrale improprio ordinario diverge. Non sono la stessa nozione.
\end{errorefrequente}
\begin{controllofinale}{impropri}
Dichiara separatamente convergenza e valore. Un'antiderivata formale non sostituisce l'esistenza dei limiti impropri.
\end{controllofinale}

Un integrale improprio \textbf{converge} soltanto quando tutti i limiti che lo definiscono esistono come numeri reali finiti. Un limite uguale a \(+\infty\) o \(-\infty\), oppure inesistente, determina la divergenza dell'integrale improprio ordinario.

\section{Definizioni e valore principale di Cauchy}\label{definizioni-e-valore-principale-di-cauchy}

Si assume \(f\in\mathcal R_{\mathrm{loc}}\) fuori dai punti singolari, cioè Riemann-integrabile su ogni intervallo compatto che non contiene singolarità.

Estremi infiniti:

\[
\displaystyle
\int_a^{+\infty}f(x)\,dx
=
\lim_{R\to+\infty}\int_a^R f(x)\,dx,
\]

\[
\displaystyle
\int_{-\infty}^{b}f(x)\,dx
=
\lim_{R\to-\infty}\int_R^b f(x)\,dx.
\]

Intera retta. Per un qualunque \(c\in\mathbb R\), la somma seguente è la definizione dell'integrale soltanto quando entrambi i termini convergono:

\[
\displaystyle
\int_{-\infty}^{+\infty}f(x)\,dx
=
\int_{-\infty}^{c}f(x)\,dx
+
\int_c^{+\infty}f(x)\,dx.
\]

\[
\displaystyle
\int_{-\infty}^{+\infty}f\ \text{converge}
\quad\Longleftrightarrow\quad
\int_{-\infty}^{c}f\ \text{e}\ \int_c^{+\infty}f\ \text{convergono}.
\]

Singolarità agli estremi:

\[
\displaystyle
\int_a^b f(x)\,dx
=
\lim_{\varepsilon\to0^+}
\int_{a+\varepsilon}^{b}f(x)\,dx,
\]

\[
\displaystyle
\int_a^b f(x)\,dx
=
\lim_{\varepsilon\to0^+}
\int_a^{b-\varepsilon}f(x)\,dx.
\]

Se entrambi gli estremi sono singolari, si sceglie un punto \(c\in(a,b)\) e si definisce

\[
\int_a^b f
=
\int_a^c f+\int_c^b f,
\]

richiedendo separatamente la convergenza finita dei due integrali.

Singolarità interna \(c\in(a,b)\):

\[
\displaystyle
\int_a^b f
=
\lim_{t\to c^-}\int_a^t f(x)\,dx
+
\lim_{t\to c^+}\int_t^b f(x)\,dx,
\]

\[
\displaystyle
\int_a^b f\ \text{converge}
\quad\Longleftrightarrow\quad
\text{entrambi i limiti sono finiti}.
\]

Valore principale di Cauchy per una singolarità interna \(c\in(a,b)\):

\[
\displaystyle
\operatorname{PV}\int_a^b f(x)\,dx
=
\lim_{\varepsilon\to0^+}
\left[
\int_a^{c-\varepsilon}f(x)\,dx
+
\int_{c+\varepsilon}^{b}f(x)\,dx
\right].
\]

Le due parti possono divergere separatamente e cancellarsi nel limite simmetrico; per questo l'esistenza del valore principale non implica la convergenza dell'integrale improprio ordinario.

Nel caso simmetrico rispetto all'origine, con \(a>0\):

\[
\displaystyle
\operatorname{PV}\int_{-a}^{a}f(x)\,dx
=
\lim_{\varepsilon\to0^+}
\left(
\int_{-a}^{-\varepsilon}f(x)\,dx
+
\int_{\varepsilon}^{a}f(x)\,dx
\right),
\]

\[
\displaystyle
\operatorname{PV}\int_{-\infty}^{+\infty}f(x)\,dx
=
\lim_{R\to+\infty}\int_{-R}^{R}f(x)\,dx.
\]

Caso particolare:

\[
\displaystyle
\operatorname{PV}\int_{-1}^{1}\frac{dx}{x}=0,
\qquad
\int_{-1}^{1}\frac{dx}{x}\ \text{diverge}.
\]

Teoria e casi multipli: \href{https://ingegnerismo.it/matematica/integrali-impropri/}{integrali impropri}.

\section{Integrali impropri campione}\label{integrali-impropri-campione}

\[
\displaystyle
\int_1^{+\infty}\frac{dx}{x^p}
=
\begin{cases}
\dfrac1{p-1}, & p>1,\\[4pt]
\text{diverge}, & p\le1,
\end{cases}
\]

\[
\displaystyle
\int_0^1\frac{dx}{x^p}
=
\begin{cases}
\dfrac1{1-p}, & p<1,\\[4pt]
\text{diverge}, & p\ge1.
\end{cases}
\]

Secondo campione fondamentale:

\[
\int_e^{+\infty}\frac{dx}{x(\ln x)^q}
=
\frac1{q-1},
\qquad q>1,
\]

\[
\int_e^{+\infty}\frac{dx}{x(\ln x)^q}\text{ diverge}
\Longleftrightarrow q\le1.
\]

Regole asintotiche pratiche per integrande positive, con \(A\) sufficientemente grande e \(\delta>0\) sufficientemente piccolo:

\[
f(x)\sim\frac{C}{x^p}\quad(x\to+\infty),\ C>0
\Longrightarrow
\int_A^{+\infty} f(x)\,dx\text{ converge}\Longleftrightarrow p>1,
\]

\[
f(x)\sim\frac{C}{(x-a)^p}\quad(x\to a^+),\ C>0
\Longrightarrow
\int_a^{a+\delta} f(x)\,dx\text{ converge}\Longleftrightarrow p<1.
\]

\section{Criteri di convergenza}\label{criteri-di-convergenza}

Criterio di Cauchy su \([a,+\infty)\):

\[
\displaystyle
\int_a^{+\infty}f\ \text{converge}
\quad\Longleftrightarrow\quad
\forall\varepsilon>0\ \exists R>a\ \forall A,B:
\ B>A>R
\Longrightarrow
\left\lvert\int_A^B f(x)\,dx\right\rvert<\varepsilon.
\]

Confronto, per \(0\le f\le g\) definitivamente:

\[
\displaystyle
\int g\ \text{converge}
\Longrightarrow
\int f\ \text{converge},
\]

\[
\displaystyle
\int f\ \text{diverge}
\Longrightarrow
\int g\ \text{diverge}.
\]

Confronto asintotico, per \(f,g>0\) definitivamente:

\[
\displaystyle
\lim\frac{f}{g}=\ell,\quad 0<\ell<+\infty
\quad\Longrightarrow\quad
\int f\ \text{e}\ \int g\ \text{hanno lo stesso carattere}.
\]

Casi unilaterali:

\[
\displaystyle
\frac fg\to0,\quad \int g\ \text{converge}
\quad\Longrightarrow\quad
\int f\ \text{converge},
\]

\[
\displaystyle
\frac fg\to+\infty,\quad \int g\ \text{diverge}
\quad\Longrightarrow\quad
\int f\ \text{diverge}.
\]

Convergenza assoluta:

\[
\displaystyle
\int\lvert f\rvert<+\infty
\quad\Longrightarrow\quad
\int f\ \text{converge},
\qquad
\left\lvert\int f\right\rvert\le\int\lvert f\rvert.
\]

Dirichlet, su \([a,+\infty)\):

\[
\displaystyle
F(R)=\int_a^R f(x)\,dx,\quad
\sup_{R\ge a}\lvert F(R)\rvert<+\infty,\quad
g\ \text{monotona},\quad g(x)\to0
\]

\[
\displaystyle
\Longrightarrow
\int_a^{+\infty}f(x)g(x)\,dx\ \text{converge}.
\]

Abel, su \([a,+\infty)\):

\[
\displaystyle
\int_a^{+\infty}f(x)\,dx\ \text{converge},\quad
g\ \text{monotona e limitata}
\]

\[
\displaystyle
\Longrightarrow
\int_a^{+\infty}f(x)g(x)\,dx\ \text{converge}.
\]

Teoria: \href{https://ingegnerismo.it/matematica/criteri-di-integrabilita/}{criteri di integrabilità}, \href{https://ingegnerismo.it/matematica/integrali-impropri/}{integrali impropri} e \href{https://ingegnerismo.it/matematica/criteri-di-abel-e-dirichlet/}{criteri di Abel e Dirichlet}.
